% §2 Generative model of sequences.  Delivers, in build order, everything a reader needs
% in order to READ thm:main-sim and nothing more: tree, GTR, clock, similarity, CFBM,
% Laplacian and Fiedler vector, eta, rho, thm:stdr-split, Pi(.), Condition C1.
% This section depends on nothing below it.
\section{Generative model of sequences}\label{sec:model}

This section fixes the generative model in the order in which it is built: the
tree, the substitution process on its edges, the clock that turns branch lengths
into times, the similarity matrix that the process induces on the leaves, and the
block structure that the tree leaves behind in that matrix. The section closes with
the spectral recovery result that makes the similarity matrix the object of
interest.

\subsection{Tree topology}\label{sec:model-topology}

Let $\mathcal{T} = (V, E)$ be a rooted latent binary phylogenetic tree whose vertex
set $V = \mathcal{X} \cup \mathcal{H}$ splits into $m$ observed terminal leaves
$\mathcal{X} = \{x_1, \dots, x_m\}$ and $m-1$ unobserved internal ancestral nodes
$\mathcal{H} = \{h_1, \dots, h_{m-1}\}$, among them a distinguished root
$r \in \mathcal{H}$ \cite{aizenbud2023spectral}. In the phylogenetic reading the
leaves are extant organisms and the internal nodes their extinct ancestors.

Every edge $(h, h') \in E$ carries a \emph{branch length} $\tau_{hh'} \ge 0$, a
measure of the distance between the two adjacent nodes it joins. We record branch
lengths through absolute node heights: each node $h$ is assigned a height
$\tau_h \ge 0$ measured upward from the leaves, so that
$\tau_{hh'} = |\tau_h - \tau_{h'}|$ and a node deeper in the past sits higher. For
a pair of leaves $i, j$ we write
\begin{equation}\label{eq:mrca_height}
   \tau_{ij} := \tau_{\mathrm{MRCA}(i,j)},
   \qquad
   d_{ij} := \sum_{e \in \mathrm{path}(i,j)} \tau_e ,
\end{equation}
for the height of their most recent common ancestor and for the total length of the
path joining them. Closely related leaves have a low MRCA and a short path; the
root is the highest node in the tree. \Cref{fig:topology-heights} draws the
convention.


The two edges descending from $r$ define the \emph{primary split} of the tree: the
leaf set $\mathcal{X}$ is partitioned into the two subtrees hanging off the root.
This single bipartition $\mathcal{X} = A \cup B$, rather than the full branching
pattern of $\mathcal{T}$, is the object we set out to recover; resolving it is the
top-down step from which the remaining tree structure follows recursively. We write
$a$ and $b$ for the sizes of the two sides, ordered so that $a \le b$ and
$a + b = m$, and we order the leaves along the split, leaves $1$ through $a$ on the
smaller side, so that each side is a contiguous block of indices.

The two sides of a split have a standard name.

\begin{definition}[Clan; adjacent clans]\label{def:clan}
A \emph{clan} of $\mathcal{T}$ is the set of leaves on one side of an edge of the
underlying unrooted tree \cite{wilkinson2007clades, semple2003phylogenetics}. Two
disjoint clans are \emph{adjacent} if their union is again a clan
\cite{aizenbud2023spectral}, equivalently if the root nodes of the two subtrees
share a neighbour.
\end{definition}

\noindent \Cref{fig:topology-clans} shows one such cut and the clan it leaves on
each side.

\begin{figure}[t]
\centering
\footnotesize
\tikzset{
  inode/.style={circle, draw, fill=black!8, minimum size=6.6mm, inner sep=0pt,
                font=\tiny},
  lfN/.style={circle, draw, fill=black!8, minimum size=6.6mm, inner sep=0pt,
              font=\tiny},
  % endpoint of a highlighted path: contour in that path's colour
  lfblue/.style={lfN, draw=blue!70!black, line width=1.1pt},
  lforng/.style={lfN, draw=orange!85!black, line width=1.1pt},
  lfA/.style={circle, draw, fill=blue!22, minimum size=6.6mm, inner sep=0pt,
              font=\tiny},
  lfB/.style={circle, draw, fill=red!22, minimum size=6.6mm, inner sep=0pt,
              font=\tiny},
  br/.style={thick, black!65},
  hl/.style={line width=1.6pt, blue!70!black},
  hlr/.style={line width=1.6pt, orange!85!black},
  guide/.style={densely dotted, black!45},
  cutB/.style={line width=1pt, green!45!black, dash pattern=on 3pt off 2pt},
}

%----------------------------------------------------------------- panel (a)
\begin{subfigure}[t]{0.49\textwidth}
\centering
\begin{tikzpicture}[x=1cm, y=1.15cm]
  % leaves -- neutral fill: this panel is about heights, not clans
  \node[lfblue] (x1) at (-2.35,0) {$x_1$};
  \node[lfblue] (x2) at (-1.55,0) {$x_2$};
  \node[lforng] (x3) at (-0.75,0) {$x_3$};
  \node[lfN]    (x4) at ( 0.30,0) {$x_4$};
  \node[lforng] (x5) at ( 1.10,0) {$x_5$};
  \node[lfN] (x6) at ( 1.90,0) {$x_6$};
  \node[lfN] (x7) at ( 2.70,0) {$x_7$};
  % internal; the two MRCAs carry the colour of their path
  \node[inode, fill=blue!30]   (u)  at (-1.95,0.80) {};
  \node[inode]                 (hA) at (-1.35,1.80) {$h_A$};
  \node[inode]                 (w)  at ( 0.70,0.80) {};
  \node[inode]                 (z)  at ( 2.30,0.80) {};
  \node[inode]                 (hB) at ( 1.50,1.35) {$h_B$};
  \node[inode, fill=orange!35] (r)  at ( 0.10,2.45) {$r$};
  % branches
  \draw[br] (r)--(hA); \draw[br] (r)--(hB);
  \draw[br] (hA)--(u); \draw[br] (hA)--(x3);
  \draw[br] (u)--(x1); \draw[br] (u)--(x2);
  \draw[br] (hB)--(w); \draw[br] (hB)--(z);
  \draw[br] (w)--(x4); \draw[br] (w)--(x5);
  \draw[br] (z)--(x6); \draw[br] (z)--(x7);
  % two edge-disjoint paths: within-clan x1 -- u -- x2, cross-clan x3 -- r -- x5
  \draw[hlr] (x3)--(hA)--(r)--(hB)--(w)--(x5);
  \draw[hl]  (x1)--(u)--(x2);
  % height axis
  \draw[black!55, -{Latex[length=1.4mm]}] (-3.15,-0.05) -- (-3.15,2.85)
        node[above, font=\scriptsize, black!55] {$\tau$};
  \foreach \y/\lbl in {0/$0$, 0.80/$\tau_{12}$, 2.45/$\tau_r$}{
     \draw[black!55] (-3.25,\y) -- (-3.05,\y);
     \node[left=0pt, font=\scriptsize, black!70] at (-3.25,\y) {\lbl};
  }
  \draw[guide] (-3.05,0.80) -- (u);
  \draw[guide] (-3.05,2.45) -- (r);
  \draw[guide] (-3.05,0) -- (x1);
  % annotations
  \node[font=\scriptsize, blue!70!black] at (-1.95,-0.62) {$d_{12}=2\tau_{12}$};
  \node[font=\scriptsize, orange!75!black, align=left, anchor=west] at (0.45,2.82)
        {$\mathrm{MRCA}(x_3,x_5)=r$\\[-1pt] so $\tau_{35}=\tau_r$};
\end{tikzpicture}
\subcaption{Heights and coalescence}\label{fig:topology-heights}
\end{subfigure}
\hfill
%----------------------------------------------------------------- panel (b)
\begin{subfigure}[t]{0.49\textwidth}
\centering
\begin{tikzpicture}[x=1cm, y=1.15cm]
  \node[lfB] (x1) at (-2.35,0) {$x_1$};
  \node[lfB] (x2) at (-1.55,0) {$x_2$};
  \node[lfB] (x3) at (-0.75,0) {$x_3$};
  \node[lfB] (x4) at ( 0.30,0) {$x_4$};
  \node[lfB] (x5) at ( 1.10,0) {$x_5$};
  \node[lfA] (x6) at ( 1.90,0) {$x_6$};
  \node[lfA] (x7) at ( 2.70,0) {$x_7$};
  \node[inode] (u)  at (-1.95,0.80) {};
  \node[inode] (hA) at (-1.35,1.80) {};
  \node[inode] (w)  at ( 0.70,0.80) {};
  \node[inode] (z)  at ( 2.30,0.80) {};
  \node[inode] (hB) at ( 1.50,1.35) {};
  \node[inode] (r)  at ( 0.10,2.45) {$r$};
  \draw[br] (r)--(hA); \draw[br] (r)--(hB);
  \draw[br] (hA)--(u); \draw[br] (hA)--(x3);
  \draw[br] (u)--(x1); \draw[br] (u)--(x2);
  \draw[br] (hB)--(w); \draw[br] (hB)--(z);
  \draw[br] (w)--(x4); \draw[br] (w)--(x5);
  \draw[br] (z)--(x6); \draw[br] (z)--(x7);
  % one edge cut, above z: C1 = {x6,x7}, C2 = everything else
  \coordinate (m1) at ($(hB)!0.5!(z)$);
  \draw[cutB] ($(m1)!0.45cm!90:(z)$) -- ($(m1)!0.45cm!-90:(z)$);
  \node[font=\scriptsize, green!45!black, align=left, anchor=west] at (2.95,1.35)
        {one edge\\[-1pt] cut};
  \node[font=\scriptsize, green!45!black] at (2.30,-0.62) {$C_1$};
  \node[font=\scriptsize, green!45!black] at (-0.55,-0.62) {$C_2$};
\end{tikzpicture}
\subcaption{Clans are edge cuts}\label{fig:topology-clans}
\end{subfigure}

\caption{A tree with $m=7$ leaves. \textbf{(a)} The blue path joins a pair that
coalesces early: its two legs are equal, so $d_{12}=2\tau_{12}$, the ultrametric
identity \eqref{eq:ultrametric}. The orange path joins a pair whose most recent
common ancestor is the root, so $\tau_{35}=\tau_r$; every pair separated by the root
behaves the same way, which is what collapses a whole block of $S$ to a single value
in \Cref{prop:cfbm}. \textbf{(b)} Single edge cuts leaves a clan on each side
(\Cref{def:clan}); the cut shown gives $C_1=\{x_6,x_7\}$ and
$C_2=\{x_1,\dots,x_5\}$.}
\label{fig:tree_topology}
\end{figure}

The primary split is therefore a pair of complementary adjacent clans, $A$ and $B$.
Rootedness and clans coexist without tension, and it is worth saying why. The model
is rooted because the clock and the sampling times introduced below are statements
about time, and time has a direction. The recovered object is a clan pair because
the similarity matrix of \Cref{def:simgraph} is a function of pairwise distances
alone and is therefore invariant to where the root is placed on the tree: the
bipartitions it can express are exactly the edge cuts of the unrooted tree. A
\emph{clade} is the special case in which the cut edge is incident to the root
\cite{wilkinson2007clades}, which is the case the primary split realizes.


\subsection{The substitution process}\label{sec:model-gtr}

Sequence data arrive by running a Markov substitution process down the edges of
$\mathcal{T}$. The process is fixed by a single rate matrix.

\begin{assumption}[GTR substitution model]\label{asm:gtr}
Sequences of length $\ell$ over a $d$-character alphabet evolve along $\mathcal{T}$
under a \emph{General Time-Reversible} (GTR) rate matrix $Q \in \R^{d \times d}$
satisfying (i)~$Q_{ij}\ge0$ for $i\ne j$, (ii)~$\sum_j Q_{ij}=0$ for all $i$,
(iii)~a stationary distribution $\pi>0$ with $\pi Q=0$, and (iv)~$\mathrm{diag}(\pi)Q$
symmetric (reversibility).
\end{assumption}

\noindent Conditions (i)--(ii) make $Q$ a generator and (iii)--(iv) make the process
reversible; together they imply that all eigenvalues of $Q$ are non-positive, and
hence that $\mathrm{tr}(Q)<0$. For a pair of \emph{adjacent} nodes $h,h'$ the
substitution matrix over the branch joining them is the matrix exponential of the
generator run for the branch length,
\begin{equation}\label{eq:transition_matrix}
   P(h \mid h') = \exp\bigl(Q\,\tau_{hh'}\bigr),
   \qquad
   P(h \mid h')_{ba} = \Prob[\,h = b \mid h' = a\,].
\end{equation}
By the Markov property the transition matrix is multiplicative along a path: for
nodes joined through an intermediate node, $P(h_1 \mid h_3) = P(h_1 \mid h_2)
P(h_2 \mid h_3)$. Two nodes separated by total branch length $t$ therefore satisfy
$P = \exp(Qt)$ whether or not they are adjacent.

\subsection{The molecular clock}\label{sec:model-clock}

So far branch lengths are abstract edge weights. The clock identifies them with
elapsed time.

\begin{assumption}[Molecular clock with contemporaneous leaves]\label{asm:clock}
Branch lengths measure evolutionary time: the substitution process along an edge is
determined by $Q$ and by the time difference between its endpoints alone, with no
per-lineage rate variation. All $m$ leaves are sampled at the same time, at height
$\tau_{x} = 0$.
\end{assumption}

\noindent \Cref{asm:clock} has one consequence, and every structural property below
is drawn from it rather than assumed separately. Since all leaves sit at height $0$,
the path from a leaf $i$ up to any ancestor $h$ has length $\tau_h$; so for every
pair of leaves both legs of the path through their MRCA have the same length,
\begin{equation}\label{eq:ultrametric}
   d\bigl(i, \mathrm{MRCA}(i,j)\bigr) = d\bigl(j, \mathrm{MRCA}(i,j)\bigr) = \tau_{ij},
   \qquad\text{hence}\qquad
   d_{ij} = 2\,\tau_{ij}.
\end{equation}
A tree with this property is \emph{ultrametric}. Pairwise path length and MRCA
height carry the same information, and either may be used interchangeably; we use
the height, because it is the quantity the tree topology fixes.

\subsection{The similarity matrix}\label{sec:model-similarity}

The substitution process is not observed; what is available is a similarity between
each pair of leaves, computed from their aligned sequences.

\begin{definition}[Pairwise similarity]\label{def:similarity}
The similarity between two nodes is the geometric mean of the determinants of the
two transition matrices between them \cite{aizenbud2023spectral},
\begin{equation}\label{eq:similarity_def}
   S(h_i, h_j) := \sqrt{\det P(h_i \mid h_j)\,\det P(h_j \mid h_i)},
   \qquad S(h_i,h_i) = 1 .
\end{equation}
\end{definition}

\begin{definition}[Similarity matrix and similarity graph]\label{def:simgraph}
The \emph{similarity matrix} $S \in \R^{m \times m}$ collects the similarities over
the observed leaves, $S_{ij} := S(x_i, x_j)$, and is symmetric with unit diagonal.
It is the weighted adjacency matrix of the \emph{similarity graph} $G$: the complete
graph whose vertices are the leaves of $\mathcal{T}$ and whose edge $(x_i,x_j)$
carries weight $S_{ij}$ \cite{aizenbud2023spectral}.
\end{definition}

Under \Cref{asm:gtr} the similarity has a closed form in the tree. For a single
branch of length $\tau$, reversibility makes the two transition matrices similar
under $\mathrm{diag}(\pi)$, so they share a determinant, and
$\det\exp(Q\tau) = \exp(\mathrm{tr}(Q)\,\tau)$ gives
$S = \exp(\mathrm{tr}(Q)\,\tau)$. Because $\det$ is multiplicative and $P$ is
multiplicative along paths, so is the similarity: for any two leaves
$S_{ij} = \exp(\mathrm{tr}(Q)\,d_{ij})$. Substituting the ultrametric identity
\eqref{eq:ultrametric},
\begin{equation}\label{eq:similarity_gtr}
    S_{ij} = \exp\!\bigl(2\,\mathrm{tr}(Q)\,\tau_{ij}\bigr).
\end{equation}
Since $\mathrm{tr}(Q)<0$, the similarity in \eqref{eq:similarity_gtr} decays
exponentially in the height of the MRCA, equivalently in the evolutionary time
separating the pair: leaves whose common ancestor is recent are similar, and
similarity falls off geometrically as that ancestor recedes into the past. The
matrix $S$ therefore records the tree's coalescence heights, one entry per pair.


\subsection{The topology inside the similarity matrix}\label{sec:model-blocks}

The tree is recoverable from $S$ because clans, and only clans, leave a rank
signature in it.

\begin{lemma}[Clans are the rank-one blocks, {\cite[Lem.~3.1]{jaffe2021spectral}}]\label{lem:rank-one}
Let $A$ and $B$ partition the leaves of an unrooted binary tree $\mathcal{T}$, and
let $S(A,B)$ denote the $|A| \times |B|$ submatrix of cross-pair similarities. Then
$S(A,B)$ is rank-one if and only if $A$ and $B$ are clans of $\mathcal{T}$.
\end{lemma}

\noindent \Cref{lem:rank-one} is due to Jaffe et al.\ and is the structural result
STDR is built on \cite{jaffe2021spectral, aizenbud2023spectral}. It is the precise
sense in which the topology of
$\mathcal{T}$ is embedded in the structure of $S$: a candidate bipartition is a
split of the tree exactly when the block of similarities it cuts collapses to rank
one. Ordering the leaves along the primary split as above, the within-clan pairs
occupy the two diagonal blocks of $S$ and the cross-clan pairs its off-diagonal
block, and \Cref{lem:rank-one} says that this off-diagonal block is rank-one at the
primary split and at no bipartition that is not a split.

Under the clock the off-diagonal block is more than rank-one --- it is constant.
This is the model the paper works in.

\begin{definition}[Clan-Floor Block Model (CFBM)]\label{def:cfbm}
A symmetric matrix $S$ with unit diagonal satisfies a localized \emph{Clan-Floor
Block Model} at a bipartition $\mathcal{X} = A \cup B$ if the cross-clan
similarities are constant at a baseline $\bze \in (0,1)$ while the intra-clan
similarities clear that baseline from above by a floor,
\begin{equation}\label{eq:cfbm_floor}
   S_{ij}=\bze\quad(i\in A,\ j\in B),
   \qquad
   \bze<\Sin^{\min}\le S_{ij}\le\Smax<1\quad(i\ne j\ \text{in the same clan}),
\end{equation}
where $\Sin^{\min}$ is the weakest intra-clan similarity and
$\Smax := \max_{i \ne j} S_{ij}$. The diagonal blocks are otherwise arbitrary. We
call the special case in which the two diagonal blocks are themselves constant the
\emph{constant-block} CFBM.
\end{definition}

\noindent \Cref{def:cfbm} is the phylogenetic instance of the constant block model
of Balakrishnan et al.\ \cite{Balakrishnan2011clustering}, in the localized form
used for latent trees by \cite{aizenbud2023spectral}: two clusters with a constant
between-cluster level, and within-cluster levels constrained only from below.


Nothing further need be assumed: \Cref{asm:gtr,asm:clock} already put $S$ in this
model at the primary split.

\begin{proposition}[The population similarity is a CFBM at the primary split]\label{prop:cfbm}
Let $S$ be the population similarity matrix of \Cref{def:simgraph} under
\Cref{asm:gtr,asm:clock}, and suppose every branch of $\mathcal{T}$ has positive
length. Then $S$ satisfies \Cref{def:cfbm} at the primary split
$\mathcal{X} = A \cup B$, with
\begin{equation}\label{eq:cfbm_from_model}
   \bze = \exp\bigl(2\,\mathrm{tr}(Q)\,\tau_r\bigr),
   \qquad
   \Sin^{\min} = \exp\bigl(2\,\mathrm{tr}(Q)\,\tau^{\max}\bigr) > \bze,
\end{equation}
where $\tau_r$ is the root height and $\tau^{\max} < \tau_r$ is the greatest
within-clan MRCA height. In particular $S_{ii}=1$ and $0 < \bze \le S_{ij} \le
\Smax < 1$ for every $i \ne j$, and the block degrees of $L = D - S$ are
$D_A=\operatorname{diag}(S_A\1_a)+b\bze I_a$ and
$D_B=\operatorname{diag}(S_B\1_b)+a\bze I_b$.
\end{proposition}

\begin{proof}
Let $i \in A$ and $j \in B$. Every path from $A$ to $B$ traverses the root, so
$\mathrm{MRCA}(i,j) = r$ and $\tau_{ij} = \tau_r$ for every cross-clan pair; this is
purely topological. By \eqref{eq:similarity_gtr} all cross-clan entries then equal
the single value $\bze = \exp(2\,\mathrm{tr}(Q)\tau_r)$, which lies in $(0,1)$
because $\mathrm{tr}(Q)<0$ and $\tau_r>0$.

Now let $i \ne j$ lie in the same clan, say $A$, with clan root $h_A$. Then
$\mathrm{MRCA}(i,j)$ is a descendant of $h_A$, so
$\tau_{ij} \le \tau^{\max} \le \max(\tau_{h_A},\tau_{h_B}) < \tau_r$, the last
inequality strict because the two root branches have positive length. As
$\mathrm{tr}(Q)<0$, \eqref{eq:similarity_gtr} is strictly decreasing in $\tau_{ij}$,
so $S_{ij} \ge \exp(2\,\mathrm{tr}(Q)\tau^{\max}) = \Sin^{\min} > \bze$. In the
other direction, $\mathrm{MRCA}(i,j)$ is a strict ancestor of both $i$ and $j$ for
$i \ne j$, so $\tau_{ij}>0$ and $S_{ij}<1$; taking the maximum over the finitely
many pairs gives $\Smax<1$. Together with $\tau_{ii}=0$, giving
$S_{ii}=\exp(0)=1$, this is \eqref{eq:cfbm_floor} and the stated range. The block
degrees follow by summing the rows of $S$ within each block and adding the constant
cross-block contribution, $b\bze$ for a row of $A$ and $a\bze$ for a row of $B$.
\end{proof}


With the baseline and the floor in hand, the primary split is visible in $S$ as two
strong diagonal blocks separated by a weaker constant off-diagonal block, and in
$L = D - S$ as the corresponding block form,
\begin{equation}\label{eq:block_forms}
   S=\begin{pmatrix} S_A & \bze J_{a\times b}\\ \bze J_{b\times a} & S_B\end{pmatrix},
   \qquad
   L=\begin{pmatrix} D_A-S_A & -\bze J_{a\times b}\\ -\bze J_{b\times a} & D_B-S_B\end{pmatrix},
\end{equation}
with $J_{a\times b}$ the all-ones block and $\bze$ the cross-clan baseline. The
diagonal blocks are whatever the two subtrees make them; only the off-diagonal block
is pinned.

\subsection{Spectral quantities}\label{sec:model-spectral}

The split is read out of $S$ spectrally, through the Laplacian.

\begin{definition}[Laplacian and Fiedler vector]\label{def:fiedler}
The \emph{combinatorial Laplacian} of the similarity graph is $L := D - S$, where
$D$ is diagonal with $D_{ii} = \sum_j S_{ij}$ \cite{aizenbud2023spectral}. Write
$0 = \lambda_1 \le \lambda_2 \le \lambda_3 \le \dots \le \lambda_m$ for its
eigenvalues. The \emph{Fiedler vector} $\vv$ is a unit eigenvector belonging to the
second-smallest eigenvalue $\lambda_2$ \cite{fiedler1975property,
vonluxburg2007tutorial}.
\end{definition}

\noindent Two scalar parameters enter the sampling rate of \Cref{sec:main}. The
first is purely topological.

\begin{definition}[Imbalance ratio]\label{def:imbalance}
For clan sizes $a \le b$ with $a+b=m$, the \emph{imbalance ratio} measures how
unevenly the primary split divides the leaves,
\begin{equation}\label{eq:imbalance_ratio}
   \eta:=\frac{b}{a}\ge1 .
\end{equation}
\end{definition}

\noindent A perfectly balanced split has $\eta=1$, while caterpillar-like lineages
reach $\eta=\mathcal O(m)$; this parameter plays the role of the balanced-cluster
constant of Balakrishnan et al.\ \cite{Balakrishnan2011clustering}. The second
parameter measures how sharply intra-clan similarities exceed cross-clan ones.

\begin{definition}[Structural margin]\label{def:margin}
Write $\Sout^{\max}$ and $\Sout^{\min}$ for the largest and smallest similarity
across the primary split, and $\Sin^{\min}$, $\Smax$ as in \eqref{eq:cfbm_floor},
so that $\Smax$ is the largest off-diagonal entry of $S$ overall. Following the
range-restriction framework of \cite{Balakrishnan2011clustering}, the
\emph{structural margin} is
\begin{equation}\label{eq:margin_definition}
   \rho := (\Sin^{\min}-\Sout^{\max})-(\Sout^{\max}-\Sout^{\min}).
\end{equation}
\end{definition}

\noindent The first bracket is the separation between the weakest within-clan pair
and the strongest cross-clan pair; the second penalizes variation inside the
cross-clan block. Under \Cref{asm:clock} that variation vanishes,
$\Sout^{\max}=\Sout^{\min}=\bze$ by \Cref{prop:cfbm}, the penalty collapses, and
\begin{equation}\label{eq:margin_under_clock}
   \rho=\Sin^{\min}-\bze > 0 .
\end{equation}

\subsection{Spectral recovery of the split}\label{sec:model-recovery}

We can now say why the similarity matrix is the object to compute. Thresholding the
Fiedler vector of the similarity graph never returns a bipartition that is not a
split of the tree.

\begin{theorem}[Consistency of the Fiedler split, {\cite[Thm.~4.2]{aizenbud2023spectral}}]\label{thm:stdr-split}
Let $G$ be the similarity graph of a binary tree $\mathcal{T}$ and let $v$ be the
Fiedler vector of $G$. Partition the leaves by the sign pattern of $v$,
\begin{equation}\label{eq:sign_rule}
   C_1=\{i:\;v_i\ge0\},\qquad C_2=\{j:\;v_j<0\}.
\end{equation}
Then $C_1$ and $C_2$ are adjacent clans of $\mathcal{T}$.
\end{theorem}

\noindent \Cref{thm:stdr-split} carries no hypothesis on tree shape or on the
substitution parameters: given the exact similarity matrix, the sign rule
\eqref{eq:sign_rule} always cuts the tree along a single edge. This is what makes
top-down spectral recovery viable at all, and it is why the recursion can merge its
two subproblems --- each side is a genuine subtree. \Cref{fig:fiedler_partitions} shows
the split on three simulated trees of increasing imbalance.

We write $\Pisp{S}$ for the bipartition \eqref{eq:sign_rule} produced by the Fiedler
vector of $L=D-S$, so \Cref{thm:stdr-split} reads: $\Pisp{S}$ is a clan pair of
$\mathcal{T}$. This is the object the rest of the paper is about --- \Cref{sec:problem}
asks at what sampling rate $\Pisp{\hat S}$ still equals $\Pisp{S}$ --- and it is
well-defined whenever $\lambda_2$ is simple.

\begin{figure}[htbp]
\centering
\begin{subfigure}[t]{0.30\textwidth}
  \captionsetup{singlelinecheck=false, justification=raggedright, skip=1pt}
  \subcaption{$\eta=1.5$}\label{fig:fied-eta01}
  \includegraphics[width=\textwidth]{Fiedler_Bipartitions/fiedler_tree_partition_eta01.png}
\end{subfigure}
\hfill{\color{black!25}\rule[-5.0cm]{0.4pt}{5.0cm}}\hfill
\begin{subfigure}[t]{0.30\textwidth}
  \captionsetup{singlelinecheck=false, justification=raggedright, skip=1pt}
  \subcaption{$\eta=4.4$}\label{fig:fied-eta05}
  \includegraphics[width=\textwidth]{Fiedler_Bipartitions/fiedler_tree_partition_eta05.png}
\end{subfigure}
\hfill{\color{black!25}\rule[-5.0cm]{0.4pt}{5.0cm}}\hfill
\begin{subfigure}[t]{0.30\textwidth}
  \captionsetup{singlelinecheck=false, justification=raggedright, skip=1pt}
  \subcaption{$\eta=15.1$}\label{fig:fied-eta15}
  \includegraphics[width=\textwidth]{Fiedler_Bipartitions/fiedler_tree_partition_eta15.png}
\end{subfigure}
\caption{$\Pisp{S}$ on three trees, one per panel. Each is drawn
from a pool of simulated Kingman trees ($m=500$ leaves, JC69 sequences of length
$\ell=10^4$) binned by the imbalance its Fiedler split realises. In each panel the
tree is on top, its branches coloured by the sign of the Fiedler vector --- grey on
the one branch that spans both signs --- and below it the Fiedler value of each leaf
is plotted at that leaf's own horizontal position in the tree. The sign pattern cuts a single edge every
time, as \Cref{thm:stdr-split} requires, and the smaller the clan it cuts off, the
larger the values carried by its leaves.}
\label{fig:fiedler_partitions}
\end{figure}


What the theorem does not say is \emph{which} edge of $\mathcal{T}$ it cuts, and it
leaves $\lambda_2$ free to be degenerate, in which case $\Pisp{S}$ is not even well
defined. One quantitative condition on the tree closes both points, and it is the
only condition beyond \Cref{asm:gtr,asm:clock} that the main result carries.

\begin{condition}[Separation margin / recoverability]\label{asm:margin}
The structural margin strictly dominates the imbalance-weighted cross-clan
boundary,
\begin{equation}\label{eq:margin_assumption}
   \rho>\eta\,\Sout^{\max},\qquad\text{equivalently}\qquad
   \Sin^{\min}>(1+\eta)\,\Sout^{\max}.
\end{equation}
\end{condition}

\noindent Under \Cref{asm:clock} that range collapses to
the single value $\bze$ by \Cref{prop:cfbm}, leaving $\Sin^{\min}>(1+\eta)\bze$.

\noindent With the model, the two governing parameters and \Cref{asm:margin} in hand,
\Cref{sec:problem} fixes what is observed and \Cref{sec:main} states the rate.
\Cref{sec:bridging} then returns to the population matrix, identifies which clan pair
$\Pisp{S}$ is, and traces how the imbalance $\eta$ moves the two spectral quantities
that the constants of that rate conceal.
